\documentclass[conference,a4paper]{IEEEtran}
\IEEEoverridecommandlockouts

\usepackage[T1]{fontenc}
\usepackage{cite}
\usepackage{amsmath,amssymb,amsfonts,mathtools}
\usepackage{algorithmic}
\usepackage{graphicx}
\usepackage{textcomp}
\usepackage{xcolor}
\usepackage{flushend}
\usepackage{fancyhdr}
\usepackage{balance}
\usepackage{multirow}
\usepackage{array}
\usepackage{tabu}
\usepackage{color}
\usepackage{ragged2e}
\usepackage{microtype}
\usepackage{url}
\def\UrlBreaks{\do\/\do-\do\_}

\def\BibTeX{{\rm B\kern-.05em{\sc i\kern-.025em b}\kern-.08em
    T\kern-.1667em\lower.7ex\hbox{E}\kern-.125emX}}
 
\renewcommand\IEEEkeywordsname{Keywords}
% No en or em dashes anywhere: IEEEtran ends the Abstract and Keywords labels with ---, and
% joins a compressed citation range with ]--[. A colon and a hyphen instead, as the Word copy
% does (build_docx.py). Two-column layout only, which is the one this paper uses.
\makeatletter
\def\abstract{\normalfont\@IEEEabskeysecsize\bfseries\textit{\abstractname}:\ \relax\@IEEEgobbleleadPARNLSP}
\def\IEEEkeywords{\normalfont\@IEEEabskeysecsize\bfseries\textit{\IEEEkeywordsname}:\ \relax\@IEEEgobbleleadPARNLSP}
\makeatother
\renewcommand{\citedash}{]-[}

\setlength{\headheight}{13mm}
% Page Layout setup
\voffset = 4mm
\headsep = 16pt
\textwidth = 182mm
\textheight = 238mm
\footskip = 12.157mm
\usepackage{hyperref}

\usepackage{pgfplots}
\usepackage{pgfplotstable}
\pgfplotsset{compat=1.18}
% Every number in the text comes from the generated macros of the runs on the paper's loop.
% Generated by analysis/summarize.py. Do not edit.
% Delays in microseconds (nearest-rank medians); residual shares in percent of B_eff.
% Sources: results/raw/2026-09-29-L-publication-0af1ac4/bound-sweep, results/raw/2026-09-29-L-publication-0af1ac4/bound-sweep-qos, results/raw/2026-09-29-L-publication-0af1ac4/gap-sweep, results/raw/2026-09-29-L-publication-0af1ac4/matrix, results/raw/2026-09-29-L-random-gap-0af1ac4/random-gap, results/raw/2026-09-30-W-publication/bound-sweep, results/raw/2026-09-30-W-publication/gap-sweep, results/raw/2026-09-30-W-publication/matrix
% Code commit: 0af1ac47ec708595b91e5c15300556f2714314c1
% Resampling of the gap sweep: random.Random(20260926), 100000 draws per cell, Python 3.14.5.
\newcommand{\WakeEpollFixedOneMs}{6.8}
\newcommand{\WakeDetectEpollFixedOneMs}{6.9}
\newcommand{\WakeEpollDefectOneMs}{300.9}
\newcommand{\WakeEpollBEffOneMs}{1007.8}
\newcommand{\WakeEpollAbsResidualPctOneMs}{0.13}
\newcommand{\WakeDetectEpollDefectOneMs}{300.6}
\newcommand{\WakeDetectThresholdEpollOneMs}{201.6}
\newcommand{\WakeEpollFixedOneHundredMs}{48.3}
\newcommand{\WakeDetectEpollFixedOneHundredMs}{48.3}
\newcommand{\WakeEpollDefectOneHundredMs}{30456.9}
\newcommand{\WakeEpollBEffOneHundredMs}{100401.4}
\newcommand{\WakeEpollAbsResidualPctOneHundredMs}{0.03}
\newcommand{\WakeDetectEpollDefectOneHundredMs}{30446.0}
\newcommand{\WakeDetectThresholdEpollOneHundredMs}{20080.3}
\newcommand{\WakeEpollFixedOneMsDefaultSlack}{7.1}
\newcommand{\WakeDetectEpollFixedOneMsDefaultSlack}{6.9}
\newcommand{\WakeEpollDefectOneMsDefaultSlack}{349.1}
\newcommand{\WakeEpollBEffOneMsDefaultSlack}{1057.0}
\newcommand{\WakeEpollAbsResidualPctOneMsDefaultSlack}{0.13}
\newcommand{\WakeDetectEpollDefectOneMsDefaultSlack}{349.2}
\newcommand{\WakeDetectThresholdEpollOneMsDefaultSlack}{211.4}
\newcommand{\WakeUringFixedOneMs}{5.4}
\newcommand{\WakeDetectUringFixedOneMs}{5.5}
\newcommand{\WakeUringDefectOneMs}{299.1}
\newcommand{\WakeUringBEffOneMs}{1006.8}
\newcommand{\WakeUringAbsResidualPctOneMs}{0.09}
\newcommand{\WakeDetectUringDefectOneMs}{299.2}
\newcommand{\WakeDetectThresholdUringOneMs}{201.4}
\newcommand{\WakeUringFixedOneHundredMs}{47.7}
\newcommand{\WakeDetectUringFixedOneHundredMs}{47.8}
\newcommand{\WakeUringDefectOneHundredMs}{30255.1}
\newcommand{\WakeUringBEffOneHundredMs}{100324.0}
\newcommand{\WakeUringAbsResidualPctOneHundredMs}{0.04}
\newcommand{\WakeDetectUringDefectOneHundredMs}{30256.2}
\newcommand{\WakeDetectThresholdUringOneHundredMs}{20064.8}
\newcommand{\WakeUringFixedOneMsDefaultSlack}{5.5}
\newcommand{\WakeDetectUringFixedOneMsDefaultSlack}{6.1}
\newcommand{\WakeUringDefectOneMsDefaultSlack}{249.2}
\newcommand{\WakeUringBEffOneMsDefaultSlack}{1007.0}
\newcommand{\WakeUringAbsResidualPctOneMsDefaultSlack}{0.13}
\newcommand{\WakeDetectUringDefectOneMsDefaultSlack}{249.2}
\newcommand{\WakeDetectThresholdUringOneMsDefaultSlack}{201.4}
\newcommand{\WakeIocpFixedOneMsTimerOneMs}{7.2}
\newcommand{\WakeDetectIocpFixedOneMsTimerOneMs}{6.9}
\newcommand{\WakeIocpDefectOneMsTimerOneMs}{1037.8}
\newcommand{\WakeIocpBEffOneMsTimerOneMs}{1516.1}
\newcommand{\WakeIocpAbsResidualPctOneMsTimerOneMs}{7.67}
\newcommand{\WakeDetectIocpDefectOneMsTimerOneMs}{1014.3}
\newcommand{\WakeDetectThresholdIocpOneMsTimerOneMs}{303.2}
\newcommand{\WakeIocpFixedOneHundredMsTimerOneMs}{14.2}
\newcommand{\WakeDetectIocpFixedOneHundredMsTimerOneMs}{14.2}
\newcommand{\WakeIocpDefectOneHundredMsTimerOneMs}{30238.6}
\newcommand{\WakeIocpBEffOneHundredMsTimerOneMs}{100273.1}
\newcommand{\WakeIocpAbsResidualPctOneHundredMsTimerOneMs}{0.37}
\newcommand{\WakeDetectIocpDefectOneHundredMsTimerOneMs}{30365.0}
\newcommand{\WakeDetectThresholdIocpOneHundredMsTimerOneMs}{20054.6}
\newcommand{\WakeIocpFixedOneMsDefaultTimer}{6.6}
\newcommand{\WakeDetectIocpFixedOneMsDefaultTimer}{5.6}
\newcommand{\WakeIocpDefectOneMsDefaultTimer}{13701.3}
\newcommand{\WakeIocpBEffOneMsDefaultTimer}{15990.4}
\newcommand{\WakeIocpAbsResidualPctOneMsDefaultTimer}{2.73}
\newcommand{\WakeDetectIocpDefectOneMsDefaultTimer}{13866.5}
\newcommand{\WakeDetectThresholdIocpOneMsDefaultTimer}{3198.1}
\newcommand{\WakeIocpFixedOneHundredMsDefaultTimer}{13.6}
\newcommand{\WakeDetectIocpFixedOneHundredMsDefaultTimer}{13.9}
\newcommand{\WakeIocpDefectOneHundredMsDefaultTimer}{49830.4}
\newcommand{\WakeIocpBEffOneHundredMsDefaultTimer}{109709.7}
\newcommand{\WakeIocpAbsResidualPctOneHundredMsDefaultTimer}{1.01}
\newcommand{\WakeDetectIocpDefectOneHundredMsDefaultTimer}{50132.2}
\newcommand{\WakeDetectThresholdIocpOneHundredMsDefaultTimer}{21941.9}
\newcommand{\WakeSweepEpollBEffOneMs}{1007.8}
\newcommand{\WakeSweepEpollAbsResidualPctOneMs}{0.16}
\newcommand{\WakeSweepEpollBEffTenMs}{10332.2}
\newcommand{\WakeSweepEpollAbsResidualPctTenMs}{0.53}
\newcommand{\WakeSweepUringBEffOneMs}{1006.9}
\newcommand{\WakeSweepUringAbsResidualPctOneMs}{0.19}
\newcommand{\WakeSweepUringBEffTenMs}{10323.2}
\newcommand{\WakeSweepUringAbsResidualPctTenMs}{0.50}
\newcommand{\WakeSweepIocpBEffOneMsTimerOneMs}{1519.7}
\newcommand{\WakeSweepIocpAbsResidualPctOneMsTimerOneMs}{32.90}
\newcommand{\WakeSweepIocpBEffTenMsTimerOneMs}{10093.2}
\newcommand{\WakeSweepIocpAbsResidualPctTenMsTimerOneMs}{4.54}
\newcommand{\WakeSweepIocpBEffOneMsDefaultTimer}{15653.4}
\newcommand{\WakeSweepIocpAbsResidualPctOneMsDefaultTimer}{2.53}
\newcommand{\WakeSweepIocpBEffTenMsDefaultTimer}{15681.6}
\newcommand{\WakeSweepIocpAbsResidualPctTenMsDefaultTimer}{2.48}
\newcommand{\WakeBoundEpollFixedBlocking}{36.1}
\newcommand{\WakeBoundEpollFixedOneMs}{6.8}
\newcommand{\WakeBoundEpollFixedTenMs}{36.0}
\newcommand{\WakeBoundEpollFixedOneHundredMs}{36.2}
\newcommand{\WakeBoundUringFixedBlocking}{46.3}
\newcommand{\WakeBoundUringFixedOneMs}{5.0}
\newcommand{\WakeBoundUringFixedTenMs}{47.0}
\newcommand{\WakeBoundUringFixedOneHundredMs}{36.1}
\newcommand{\WakeBoundIocpFixedBlockingTimerOneMs}{7.1}
\newcommand{\WakeBoundIocpFixedOneMsTimerOneMs}{7.3}
\newcommand{\WakeBoundIocpFixedTenMsTimerOneMs}{8.1}
\newcommand{\WakeBoundIocpFixedOneHundredMsTimerOneMs}{8.1}
\newcommand{\WakeBoundIocpFixedBlockingDefaultTimer}{6.9}
\newcommand{\WakeBoundIocpFixedOneMsDefaultTimer}{7.1}
\newcommand{\WakeBoundIocpFixedTenMsDefaultTimer}{6.9}
\newcommand{\WakeBoundIocpFixedOneHundredMsDefaultTimer}{8.5}
\newcommand{\WakeBoundQosEpollFixedBlockingQosZeroUs}{5.1}
\newcommand{\WakeBoundQosEpollFixedOneMsQosZeroUs}{2.8}
\newcommand{\WakeBoundQosEpollFixedTenMsQosZeroUs}{5.3}
\newcommand{\WakeBoundQosEpollFixedOneHundredMsQosZeroUs}{5.3}
\newcommand{\WakeBoundQosUringFixedBlockingQosZeroUs}{5.1}
\newcommand{\WakeBoundQosUringFixedOneMsQosZeroUs}{2.9}
\newcommand{\WakeBoundQosUringFixedTenMsQosZeroUs}{5.2}
\newcommand{\WakeBoundQosUringFixedOneHundredMsQosZeroUs}{5.2}
\newcommand{\WakeBoundEpollRatioMax}{5.357}
\newcommand{\WakeBoundUringRatioMax}{9.409}
\newcommand{\WakeBoundIocpRatioMaxTimerOneMs}{1.141}
\newcommand{\WakeBoundIocpRatioMaxDefaultTimer}{1.232}
\newcommand{\WakeBoundQosEpollRatioMaxQosZeroUs}{1.924}
\newcommand{\WakeBoundQosUringRatioMaxQosZeroUs}{1.795}
% Derived: design constants, host facts, aggregates, resampling and random gaps.
\newcommand{\WakeTestPosts}{21}
\newcommand{\WakeTestDivisor}{5}
\newcommand{\WakeHOneLimitPct}{5}
\newcommand{\WakeHTwoLimit}{1.25}
\newcommand{\WakeHFiveLimit}{1.25}
\newcommand{\WakeCalibPosts}{11}
\newcommand{\WakeMatrixPosts}{51}
\newcommand{\WakeBoundPosts}{51}
\newcommand{\WakeWarmupPosts}{5}
\newcommand{\WakeMatrixGapFrac}{1.7}
\newcommand{\WakeBoundGapMs}{2}
\newcommand{\WakeSweepSteps}{30}
\newcommand{\WakeSweepStepPosts}{21}
\newcommand{\WakeSweepCellPosts}{630}
\newcommand{\WakeSweepGapLoFrac}{0.1}
\newcommand{\WakeSweepGapHiFrac}{3.0}
\newcommand{\WakeSweepPhaseLoFrac}{0.05}
\newcommand{\WakeSweepPhaseHiFrac}{2.95}
\newcommand{\WakeBSmallMs}{1}
\newcommand{\WakeBMidMs}{10}
\newcommand{\WakeBLargeMs}{100}
\newcommand{\WakeSlackNs}{1}
\newcommand{\WakeSlackDefaultUs}{50}
\newcommand{\WakeQosUs}{0}
\newcommand{\WakeTimerRequestedMs}{1}
\newcommand{\WakeWinTimerResolutionMs}{1.0}
\newcommand{\WakeRunsWindows}{136}
\newcommand{\WakeRunsLinux}{148}
\newcommand{\WakeRunsRandom}{8}
\newcommand{\WakeQuietThresholdPct}{95}
\newcommand{\WakeQuietMeanPct}{97.3}
\newcommand{\WakeHostWindowsCpu}{Ryzen 5 3600}
\newcommand{\WakeHostWindowsOs}{Windows 11}
\newcommand{\WakeHostLinuxCpu}{Ryzen 7 5800H}
\newcommand{\WakeHostLinuxKernel}{7.2.3}
\newcommand{\WakeCpuidleDriver}{acpi\_idle}
\newcommand{\WakeCpuidleGovernor}{menu}
\newcommand{\WakeCpuidleDeepest}{C3}
\newcommand{\WakeCpuidleDeepestLatencyUs}{350}
\newcommand{\WakeCompilerLinux}{clang 22.1.8}
\newcommand{\WakeCodeCommitLinux}{0af1ac4}
\newcommand{\WakeCompilerWindows}{clang-cl 22.1.0}
\newcommand{\WakeCodeCommitWindows}{0af1ac4}
\newcommand{\WakeHDroppedBoundUs}{100}
\newcommand{\WakeFixedMedianMaxUs}{48.3}
\newcommand{\WakeFigPostsShown}{300}
\newcommand{\WakeHOneLinuxMaxPct}{0.53}
\newcommand{\WakeHOneWindowsConfigs}{6}
\newcommand{\WakeHOneWindowsPassConfigs}{5}
\newcommand{\WakeHOneWindowsPassMaxPct}{4.54}
\newcommand{\WakeDetectDefectConfigs}{10}
\newcommand{\WakeDetectDefectFlagged}{10}
\newcommand{\WakeDetectBDefectFlagged}{10}
\newcommand{\WakeDetectFactorDefectMin}{1.24}
\newcommand{\WakeDetectFactorDefectMax}{4.34}
\newcommand{\WakeDetectFactorBDefectMin}{1.25}
\newcommand{\WakeDetectFactorBDefectMax}{69.3}
\newcommand{\WakeDetectConfigsLinux}{6}
\newcommand{\WakeDetectConfigsWindows}{4}
\newcommand{\WakeDetectFixedConfigs}{10}
\newcommand{\WakeDetectFixedFlagged}{0}
\newcommand{\WakeDetectBFixedFlagged}{0}
\newcommand{\WakeDetectFactorFixedMax}{0.03}
\newcommand{\WakeDetectFactorBFixedMax}{0.03}
\newcommand{\WakeSweepRunsTotal}{240}
\newcommand{\WakeSweepMissedTotal}{41}
\newcommand{\WakeSweepMissedLateTotal}{37}
\newcommand{\WakeSweepMissedBTotal}{31}
\newcommand{\WakeSweepEpollMissedOneMs}{6}
\newcommand{\WakeSweepEpollMinMedianOneMs}{74.5}
\newcommand{\WakeSweepEpollMinMedianGapMsOneMs}{2.9}
\newcommand{\WakeSweepEpollMissedTenMs}{6}
\newcommand{\WakeSweepEpollMinMedianTenMs}{39.9}
\newcommand{\WakeSweepEpollMinMedianGapMsTenMs}{20.0}
\newcommand{\WakeSweepUringMissedOneMs}{6}
\newcommand{\WakeSweepUringMinMedianOneMs}{71.4}
\newcommand{\WakeSweepUringMinMedianGapMsOneMs}{2.9}
\newcommand{\WakeSweepUringMissedTenMs}{6}
\newcommand{\WakeSweepUringMinMedianTenMs}{40.7}
\newcommand{\WakeSweepUringMinMedianGapMsTenMs}{20.0}
\newcommand{\WakeSweepIocpMissedOneMsTimerOneMs}{0}
\newcommand{\WakeSweepIocpMinMedianOneMsTimerOneMs}{994.6}
\newcommand{\WakeSweepIocpMissedTenMsTimerOneMs}{6}
\newcommand{\WakeSweepIocpMinMedianTenMsTimerOneMs}{1082.7}
\newcommand{\WakeSweepIocpMissedOneMsDefaultTimer}{6}
\newcommand{\WakeSweepIocpMinMedianOneMsDefaultTimer}{1186.8}
\newcommand{\WakeSweepIocpMissedTenMsDefaultTimer}{5}
\newcommand{\WakeSweepIocpMinMedianTenMsDefaultTimer}{1179.3}
\newcommand{\WakeResampleTrials}{100000}
\newcommand{\WakeResampleMissMaxPct}{0.08}
\newcommand{\WakeResampleMissBMaxPct}{0.08}
\newcommand{\WakeBinomialMissPct}{0.097}
\newcommand{\WakeGridUs}{500}
\newcommand{\WakeGridTolUs}{60}
\newcommand{\WakeGridUniformPct}{24}
\newcommand{\WakeGridPostsIocpOneMsTimerOneMs}{502}
\newcommand{\WakeBandPostsIocpOneMsTimerOneMs}{405}
\newcommand{\WakeBandLoUsIocpOneMsTimerOneMs}{900}
\newcommand{\WakeBandHiUsIocpOneMsTimerOneMs}{1100}
\newcommand{\WakeEpollBEffExcessOneMs}{7.8}
\newcommand{\WakeEpollBEffExcessPctOneMs}{0.8}
\newcommand{\WakeEpollBEffExcessOneHundredMs}{401.4}
\newcommand{\WakeEpollBEffExcessPctOneHundredMs}{0.4}
\newcommand{\WakeEpollBEffExcessOneMsDefaultSlack}{57.0}
\newcommand{\WakeEpollBEffExcessPctOneMsDefaultSlack}{5.7}
\newcommand{\WakeUringBEffExcessOneMs}{6.8}
\newcommand{\WakeUringBEffExcessPctOneMs}{0.7}
\newcommand{\WakeUringBEffExcessOneHundredMs}{324.0}
\newcommand{\WakeUringBEffExcessPctOneHundredMs}{0.3}
\newcommand{\WakeUringBEffExcessOneMsDefaultSlack}{7.0}
\newcommand{\WakeUringBEffExcessPctOneMsDefaultSlack}{0.7}
\newcommand{\WakeIocpBEffExcessOneMsTimerOneMs}{516.1}
\newcommand{\WakeIocpBEffExcessPctOneMsTimerOneMs}{51.6}
\newcommand{\WakeIocpBEffExcessOneHundredMsTimerOneMs}{273.1}
\newcommand{\WakeIocpBEffExcessPctOneHundredMsTimerOneMs}{0.3}
\newcommand{\WakeIocpBEffExcessOneMsDefaultTimer}{14990.4}
\newcommand{\WakeIocpBEffExcessPctOneMsDefaultTimer}{1499.0}
\newcommand{\WakeIocpBEffExcessOneHundredMsDefaultTimer}{9709.7}
\newcommand{\WakeIocpBEffExcessPctOneHundredMsDefaultTimer}{9.7}
\newcommand{\WakeSweepEpollBEffExcessOneMs}{7.8}
\newcommand{\WakeSweepEpollBEffExcessPctOneMs}{0.8}
\newcommand{\WakeSweepEpollBEffExcessTenMs}{332.2}
\newcommand{\WakeSweepEpollBEffExcessPctTenMs}{3.3}
\newcommand{\WakeSweepUringBEffExcessOneMs}{6.9}
\newcommand{\WakeSweepUringBEffExcessPctOneMs}{0.7}
\newcommand{\WakeSweepUringBEffExcessTenMs}{323.2}
\newcommand{\WakeSweepUringBEffExcessPctTenMs}{3.2}
\newcommand{\WakeSweepIocpBEffExcessOneMsTimerOneMs}{519.7}
\newcommand{\WakeSweepIocpBEffExcessPctOneMsTimerOneMs}{52.0}
\newcommand{\WakeSweepIocpBEffExcessTenMsTimerOneMs}{93.2}
\newcommand{\WakeSweepIocpBEffExcessPctTenMsTimerOneMs}{0.9}
\newcommand{\WakeSweepIocpBEffExcessOneMsDefaultTimer}{14653.4}
\newcommand{\WakeSweepIocpBEffExcessPctOneMsDefaultTimer}{1465.3}
\newcommand{\WakeSweepIocpBEffExcessTenMsDefaultTimer}{5681.6}
\newcommand{\WakeSweepIocpBEffExcessPctTenMsDefaultTimer}{56.8}
\newcommand{\WakeEpollSlackEffectOneMs}{49.1}
\newcommand{\WakeUringSlackEffectOneMs}{0.3}
\newcommand{\WakeWinDefaultTickMs}{15.7}
\newcommand{\WakeIocpTicksOneHundredMsDefaultTimer}{7}
\newcommand{\WakeIocpBEffOverTicksOneHundredMsDefaultTimer}{135.9}
\newcommand{\WakeIocpDefectMaxMsDefaultTimer}{16.4}
\newcommand{\WakeRandomTrials}{1000}
\newcommand{\WakeRandomGapLo}{0.1}
\newcommand{\WakeRandomGapHi}{3.0}
\newcommand{\WakeRandomEpollFixedFlaggedOneMs}{0}
\newcommand{\WakeRandomEpollFixedFlaggedBOneMs}{0}
\newcommand{\WakeRandomEpollDefectFlaggedOneMs}{999}
\newcommand{\WakeRandomEpollDefectFlaggedBOneMs}{999}
\newcommand{\WakeRandomEpollFixedFlaggedTenMs}{0}
\newcommand{\WakeRandomEpollFixedFlaggedBTenMs}{0}
\newcommand{\WakeRandomEpollDefectFlaggedTenMs}{999}
\newcommand{\WakeRandomEpollDefectFlaggedBTenMs}{1000}
\newcommand{\WakeRandomUringFixedFlaggedOneMs}{0}
\newcommand{\WakeRandomUringFixedFlaggedBOneMs}{0}
\newcommand{\WakeRandomUringDefectFlaggedOneMs}{998}
\newcommand{\WakeRandomUringDefectFlaggedBOneMs}{998}
\newcommand{\WakeRandomUringFixedFlaggedTenMs}{0}
\newcommand{\WakeRandomUringFixedFlaggedBTenMs}{0}
\newcommand{\WakeRandomUringDefectFlaggedTenMs}{1000}
\newcommand{\WakeRandomUringDefectFlaggedBTenMs}{1000}
\newcommand{\WakeRandomDefectTrials}{4000}
\newcommand{\WakeRandomDefectMissed}{4}
\newcommand{\WakeRandomDefectMissedB}{3}
\newcommand{\WakeRandomFixedTrials}{4000}
\newcommand{\WakeRandomFixedFlagged}{0}
\newcommand{\WakeRandomFixedFlaggedB}{0}
\newcommand{\WakeRandomDefectMissPct}{0.10}
\newcommand{\WakeRandomDefectMissBPct}{0.07}
\newcommand{\WakeRandomDefectMissUpperPct}{0.23}
\newcommand{\WakeRandomDefectMissUpperBPct}{0.19}
\newcommand{\WakeRandomAbsResidualPctMax}{0.50}
\newcommand{\WakeRandomFixedFlaggedUpperPct}{0.07}

\newcommand{\Beff}{B_{\mathrm{eff}}}
% Code availability. The repository is published, so the phrase is "is available".
% build_docx.py reads the same macro, so the PDF and the Word copy always agree.
\newcommand{\PaperRepoAvailability}{is available}
\newcommand{\PaperRepo}{\url{https://github.com/Alex-Tsvetanov/paper-wake-defect}}
% Page ranges in refs.bib: BibTeX dashifies the argument, which is dropped; one hyphen prints.
\newcommand{\BibHyphen}[1]{-}
% A line break in the Word title only (the template's larger title font wraps one word onto
% the second line); nothing in LaTeX. build_docx.py turns it into a break.
\newcommand{\WordTitleBreak}{}

\begin{document}
\bstctlcite{IEEEexample:BSTcontrol}

\title{{\fontsize{21}{14.4}\selectfont
Bounded-Wait Wake Defects \WordTitleBreak in Portable Event Loops}%
\thanks{\textsuperscript{1}Alex I. Tsvetanov is with the Faculty of Computer Systems and Technologies, Technical University of Sofia, Sofia, Bulgaria (e-mail: \href{mailto:atsvetanov@tu-sofia.bg}{atsvetanov@tu-sofia.bg}).}%
}

\author{
\IEEEauthorblockN{{\fontsize{14}{14.4}\selectfont
Alex I. Tsvetanov\textsuperscript{1}}}
}

\maketitle

\pagenumbering{arabic}
\setcounter{page}{1}
\renewcommand{\headrulewidth}{0pt}
\renewcommand{\footrulewidth}{0pt}
\thispagestyle{fancy}
% The running header, as in template.docx (straight quotes under T1).
	\chead{{\fontsize{10}{13.0}\selectfont
34th National Conference with International Participation "Telecom 2026", November 26 - 28, 2026, Sofia, Bulgaria}  }
    \fancyfoot[C]{\thepage}

\begin{abstract}
An idle event-loop worker parks in a wait with a timeout. When another thread posts a task,
the worker must be woken. If the wake is missing, or signals an object the wait does not
watch, the task still runs, but only when the timeout expires, so tests of eventual delivery
pass. Such missed wakes have been reported in widely used event loops and runtimes. We describe
the defect for epoll, io\_uring and I/O completion ports, model the delay it causes, and
propose a portable test that detects it from latency alone. We measure both on a minimal event
loop written for this study, whose correct and defective builds differ in one compile-time
switch that guards one site per backend. On Linux the median absolute error of the model is at
most \WakeHOneLinuxMaxPct\,\% of the wait period. On Windows the model holds with the default
timer and fails for a \WakeBSmallMs~ms bound under a \WakeTimerRequestedMs~ms timer. At a fixed
gap the test flagged \WakeDetectDefectFlagged\ of \WakeDetectDefectConfigs\ defective and
\WakeDetectFixedFlagged\ of \WakeDetectFixedConfigs\ correct configurations, but a gap just
below a multiple of the wait period hides the defect. With random gaps, in an exploratory run,
it missed \WakeRandomDefectMissed\ of \WakeRandomDefectTrials\ defective trials and flagged
\WakeRandomFixedFlagged\ of \WakeRandomFixedTrials\ correct ones. With a correct wake, the
delay on an idle Linux core still depends on the wait bound, also under a \WakeQosUs~$\mu$s CPU
latency request.
\end{abstract}

\begin{IEEEkeywords}
event loop, lost wake-up, epoll, io\_uring, I/O completion ports, timer resolution, processor
idle states, performance testing
\end{IEEEkeywords}

\section{Introduction}

Servers that multiplex many connections over a few threads run an event loop on each worker
thread. A worker waits for I/O events and, when nothing arrives, sleeps in the kernel. Work
also comes from other threads: a timer fires, another worker hands over a task, or the
application posts a callback. The posting thread must then interrupt the wait.

Many loops bound the wait with a timeout, to check a stop flag or to service timers. The bound
turns a missing wake-up into a delay instead of a hang: the task runs when the wait returns on
its own. A functional test that checks that a posted task eventually runs therefore passes. The
defect shows as latency, and mainly on an idle worker, because a busy one leaves its wait for
other reasons and picks the task up on the way.

Studies of real bugs suggest why such defects occur and go unnoticed. In a study of concurrency
bugs, atomicity and order violations make up almost all non-deadlock bugs, and one example has
the shape of a lost wake-up: a completion callback clears a pending flag before the waiter sets
it, and the waiter hangs~\cite{lu2008learning}. In a study of timeout bugs in cloud server
systems, most produced no error message, and some degraded performance, for example through
idle time behind a timeout that was too long~\cite{dai2018understanding}. Missed wake-ups of
event loops occur in widely used software. In Netty, a task submitted from another thread
could skip the selector wake-up and wait until the select timeout
expired~\cite{lee2016nettywakeup}. A race between \texttt{epoll\_wait} and \texttt{epoll\_ctl}
in Linux lost the \texttt{eventfd} event with which Boost.Asio wakes its loop for a handler
posted from another thread; in the reporter's test program, \texttt{epoll\_wait} returned only
at its timeout~\cite{neunhoffer2019epollbug,neunhoffer2020asio,penyaev2020epollfix}. Related
defects caused hangs. In libevent, the loop was not woken when another thread deleted its last
pending event~\cite{khuzhin2018libevent}. A lost wake-up in the Tokio scheduler let a worker
park while a spawned task was queued~\cite{lerche2019tokio}, and the .NET thread pool could
leave a queued work item unprocessed unless more work was queued~\cite{sadov2025threadpool}.
Conversely, Asio, when built for and running on Windows versions before Vista, bounds each wait
for completions, because the call can appear stuck while completions are
queued~\cite{kohlhoff2026asioiocp}: a bounded wait as a safety net.

We call this a \emph{bounded-wait wake defect}. We describe how it arises with the kernel
interfaces used by portable C++ servers (Section~\ref{sec:background}), model its delay
(Section~\ref{sec:model}), and propose a portable test that needs no access to the internals
of the loop (Section~\ref{sec:test}). We measure both on a minimal event loop written for this
study, with epoll, io\_uring and IOCP backends and a compile-time switch that seeds the defect
in each (Sections~\ref{sec:method} and~\ref{sec:results}). The repository of the
study, with the loop, its tests, the probe, the run scripts, the raw runs, the sanitizer records and
the hypotheses with their revision log, \PaperRepoAvailability\ at \PaperRepo.

\section{Background}\label{sec:background}

\emph{epoll.} The worker waits on a set of file descriptors in \texttt{epoll\_wait}. The usual
cross-thread wake is an \texttt{eventfd}~\cite{manpages2026eventfd} in the same set: the poster
writes to it and the wait returns.

\emph{io\_uring.} The worker waits for completions in
\texttt{io\_\allowbreak{}uring\_\allowbreak{}enter}~\cite{liburing2026iouringenter}.
A write to an \texttt{eventfd} ends the wait only if the ring watches that descriptor, for
example with a pending
\texttt{IORING\_\allowbreak{}OP\_\allowbreak{}POLL\_\allowbreak{}ADD}.

\emph{I/O completion ports.} On Windows the worker waits in
\texttt{GetQueued\allowbreak{}Completion\allowbreak{}Status}, and another thread wakes it by
queueing a packet with
\texttt{PostQueued\allowbreak{}Completion\allowbreak{}Status}~\cite{microsoft2026iocp}.

A portable loop wraps these behind one post operation. It must ensure that the wait observes
every posted task, even when several posts share one signal. The defect class is a failure of
that guarantee.

\section{Delay model}\label{sec:model}

Consider one worker that alternates between a wait bounded by $B$ and running the tasks it
finds. While idle, it returns from the wait with a period $\Beff \ge B$. On Linux, $\Beff$
exceeds $B$ by a small measured amount, and the timer slack of the thread can add to
it~\cite{manpages2026timerslack}. On Windows, the wait ends on a tick of the timer resolution in
effect for the process~\cite{microsoft2026timebeginperiod}, so $\Beff$ can be much larger than
$B$. Let the phase $\phi$ of a post be the time since the worker last parked. If the post does
not end the wait, the task runs at the next timed return, after
\begin{equation}
  d(\phi) = \Beff - (\phi \bmod \Beff).
  \label{eq:model}
\end{equation}
With a working wake, the delay is the cost of waking a parked thread and does not follow
(\ref{eq:model}). At uniformly random phase, the defect adds half a period on average to the
latency of a task posted to an idle worker. Where timers are routed through the post path,
every timer is delayed the same way.

\section{Detection test}\label{sec:test}

The test needs only the public post operation, a clock and the configured bound. It lets a
one-worker loop go idle, then posts $N = \WakeTestPosts$ tasks, each after the previous one has
run and a gap has passed. It fails when the median delay exceeds $\Beff/\WakeTestDivisor$.
Waiting for each task keeps the worker idle, so that no other event wakes it and hides the
defect. The median ignores occasional slow posts, such as one that wakes a core from a deeper
idle state~\cite{schone2015wakeup}, and scheduler noise.

The gap sets the phase. With a fixed gap all posts share one phase, and a gap that ends in the
last $\Beff/\WakeTestDivisor$ of a period lets a defective loop answer within
$\Beff/\WakeTestDivisor$: the test misses. With a random gap the phase is close to uniform, and
a defective loop is missed only if more than half of the posts land in that part of the period,
which for $N = \WakeTestPosts$ has probability \WakeBinomialMissPct\,\%. $\Beff$ is known only
from a defective build; a test of an unknown build can use $B/\WakeTestDivisor$, which is no
larger.

The test is a latency oracle. For an omitted or misdirected wake, as in our seeded builds, race
detectors have nothing to find: no access is racy, and
ThreadSanitizer~\cite{serebryany2009threadsanitizer} reports nothing on our defective builds.
Schedule exploration in the style of CHESS~\cite{musuvathi2008finding} would not report it
either, since every schedule eventually delivers the task. To measure the power of the test we
seed the defect with one compile-time switch that guards one site per backend, where it removes
or misdirects the wake. This is fault seeding, as in mutation testing~\cite{jia2011analysis},
with hand-written faults. Related work on wake-up cost in Linux event
notification~\cite{nolte2023hawen}, on the exit latency of idle states~\cite{schone2015wakeup}
and on tail latency from power saving~\cite{li2014tales} measures the cost of a wake that
happens; our test detects one that does not.

\section{Method}\label{sec:method}

\emph{Loop and arms.} The loop has one worker thread and no I/O, timers or second worker. A
post pushes its task onto a lock-free stack with a compare-and-swap and wakes the worker only if
the stack was empty. After every return from its wait, the worker takes the whole stack with
one exchange and runs the tasks in post order. The epoll backend writes an \texttt{eventfd}
registered edge-triggered and waits in \texttt{epoll\_\allowbreak{}pwait2} with the bound $B$.
The io\_uring backend writes an \texttt{eventfd} that a multishot poll in the worker's ring
watches, and waits in \texttt{io\_\allowbreak{}uring\_\allowbreak{}enter} with a timeout
argument. The IOCP backend posts a completion packet and waits in
\texttt{GetQueued\allowbreak{}Completion\allowbreak{}Status} with $B$ rounded up to whole
milliseconds. $B$ is a run-time parameter, and the wait can also block. Stop has its own path
in every backend. The defect switch guards one site per backend, and a test checks that there
is no other: epoll and IOCP omit the wake, and io\_uring writes its \texttt{eventfd} but never
arms the poll that watches it.

\emph{Probe.} A standalone program starts the loop, measures $\Beff$ in the defective build
from \WakeCalibPosts\ calibration posts, and after a few warm-up posts records the phase, delay
and prediction of each post. Each post's residual is taken against the $\Beff$ of its own run.
The gap sweep runs the defective build at $B = \WakeBSmallMs$ and \WakeBMidMs~ms, with
\WakeSweepSteps\ runs of \WakeSweepStepPosts\ posts per configuration, at gaps of
$\WakeSweepGapLoFrac\,B$ to $\WakeSweepGapHiFrac\,B$ on Linux and target phases of
$\WakeSweepPhaseLoFrac\,\Beff$ to $\WakeSweepPhaseHiFrac\,\Beff$ on Windows. The bound sweep
runs the correct build with \WakeBoundPosts\ posts at a \WakeBoundGapMs~ms gap, for
$B \in \{\WakeBSmallMs, \WakeBMidMs, \WakeBLargeMs\}$~ms and a blocking wait; on Linux it is
repeated under a CPU latency request of \WakeQosUs~$\mu$s on
\texttt{/dev/cpu\_dma\_latency}~\cite{kernel2026cpuidle}. The detection matrix runs both builds
at $B = \WakeBSmallMs$ and \WakeBLargeMs~ms, with \WakeMatrixPosts\ posts at a gap of
$\WakeMatrixGapFrac\,B$, whose first \WakeTestPosts\ form the test. An exploratory random-gap
design on Linux runs both builds of epoll and io\_uring at $B = \WakeBSmallMs$ and
\WakeBMidMs~ms, \WakeRandomTrials\ trials of \WakeTestPosts\ posts per cell, with gaps uniform
in $[\WakeRandomGapLo\,B, \WakeRandomGapHi\,B]$ and a recorded seed. Windows has no such run;
for both systems we instead resample the gap sweep, \WakeResampleTrials\ draws of
\WakeTestPosts\ posts without replacement per configuration, with a recorded seed.

\emph{Hypotheses.} The hypotheses and their acceptance criteria were fixed before any
publication run (Table~\ref{tab:hyp}), as their dated revision log in the repository records.
Revisions after smoke runs, before any publication run, renumbered and reworded them as H1 to
H4 (a prediction from the measured phase and $\Beff$, a fixed-gap test of independence from
$B$, a threshold of $\Beff/\WakeTestDivisor$ instead of $B/\WakeTestDivisor$, an explicit
throughput criterion). They also added the exploratory H5 after the smoke data contradicted H2,
which was kept as written. The first version's bound of \WakeHDroppedBoundUs~$\mu$s on the
correct loop's median and its threshold $B/\WakeTestDivisor$ would change no verdict (largest
correct median \WakeFixedMedianMaxUs~$\mu$s; the same H3 verdicts). H4 was then deferred and not
run for this paper, and we make no claim about it. The hypotheses were first measured on an
earlier implementation, which we withdrew; none of its data is used. After those runs and a
review, the random-gap design and the resampling of the gap sweep were added as exploratory.
Before any sanitizer record or run on the present loop, a last revision made the present loop
the measured implementation: the defect is the loop built with the switch on, the fix is the
loop with its wake compiled in, and hypotheses, thresholds and designs are unchanged.

\begin{table}[t]
\caption{Pre-specified hypotheses and their outcomes}
\label{tab:hyp}
\centering
\footnotesize
\renewcommand{\arraystretch}{1.1}
\begin{tabular}{@{}>{\raggedright\arraybackslash}p{0.45cm}>{\raggedright\arraybackslash}p{3.3cm}>{\raggedright\arraybackslash}p{4.25cm}@{}}
\hline
 & Criterion & Outcome \\
\hline
H1 & Defective loop: median $|$residual$|$ of (\ref{eq:model}) below \WakeHOneLimitPct\,\% of $\Beff$
   & Pass on Linux (largest \WakeHOneLinuxMaxPct\,\%) and on Windows with the default timer
     (largest \WakeIocpAbsResidualPctOneMsDefaultTimer\,\%); fail on Windows with the
     \WakeTimerRequestedMs~ms timer at $B = \WakeBSmallMs$~ms
     (\WakeSweepIocpAbsResidualPctOneMsTimerOneMs\,\% in the sweep,
     \WakeIocpAbsResidualPctOneMsTimerOneMs\,\% in the matrix) \\
H2 & Correct loop: largest to smallest median over $B$ at most \WakeHTwoLimit
   & Fail on Linux (epoll \WakeBoundEpollRatioMax, io\_uring \WakeBoundUringRatioMax); pass on
     Windows (\WakeBoundIocpRatioMaxTimerOneMs\ with the \WakeTimerRequestedMs~ms timer,
     \WakeBoundIocpRatioMaxDefaultTimer\ with the default) \\
H3 & The \WakeTestPosts-post test flags every defective and no correct configuration
   & Pass at a gap of $\WakeMatrixGapFrac\,B$: \WakeDetectDefectFlagged\ of \WakeDetectDefectConfigs\
     defective and \WakeDetectFixedFlagged\ of \WakeDetectFixedConfigs\ correct flagged \\
H4 & No measurable throughput cost under load & Deferred, not run; no claim \\
H5 & Exploratory: a \WakeQosUs~$\mu$s latency request brings the H2 ratio to \WakeHFiveLimit\ or below
   & Fail for epoll (\WakeBoundQosEpollRatioMaxQosZeroUs) and io\_uring
     (\WakeBoundQosUringRatioMaxQosZeroUs) \\
\hline
\end{tabular}
\end{table}

\emph{Machines.} Linux runs use a \WakeHostLinuxCpu\ laptop (boost off, performance governor,
kernel \WakeHostLinuxKernel) with the \WakeCpuidleDriver\ idle driver and the
\WakeCpuidleGovernor\ governor~\cite{kernel2026cpuidle}; its deepest idle state,
\WakeCpuidleDeepest, reports an exit latency of \WakeCpuidleDeepestLatencyUs~$\mu$s. Windows
runs use a \WakeHostWindowsCpu\ desktop with \WakeHostWindowsOs\ and the power plan recorded
with the run, whose minimum processor state on mains power is set to its highest value (checked
on the host after the runs). The runs measure the loop at commit
\texttt{\WakeCodeCommitLinux}, built with \WakeCompilerLinux\ and \WakeCompilerWindows; the
compiled files' hashes, recorded with each run, match that commit. The timer slack is \WakeSlackNs~ns, except in the $B = \WakeBSmallMs$~ms
matrix cells repeated at the default of \WakeSlackDefaultUs~$\mu$s. Windows runs use the default
timer resolution or \WakeTimerRequestedMs~ms from \texttt{timeBeginPeriod} (throttling opted
out) and time gaps with a high-resolution waitable timer. The cell order is shuffled with a
recorded seed, and the Windows run starts only when the processor is on average at least
\WakeQuietThresholdPct\,\% idle.

\emph{Validity.} A run counts only if every first-party file compiled into its binaries also
passed the loop's tests and a reduced run of every design, with the same configuration and
compiler. The sanitizers were AddressSanitizer with UndefinedBehaviorSanitizer,
ThreadSanitizer and MemorySanitizer on Linux, and AddressSanitizer with clang-cl on Windows.
None of them reported anything. On both systems the hashes of the compiled files are checked
against these records before any cell runs. These records, and those made later for the published
commit, are in the repository. The tests pin the loop's semantics in both builds
and include wake detectors, which pass in every correct build and detect the seeded defect in
every defective build.

\section{Results}\label{sec:results}

\emph{Delay model.} Fig.~\ref{fig:sawtooth} shows delay against phase at $B = \WakeBSmallMs$~ms.
On Linux the defective loops follow (\ref{eq:model}): the median absolute residual is at most
\WakeHOneLinuxMaxPct\,\% of $\Beff$ over the gap sweeps and matrix cells, and
\WakeRandomAbsResidualPctMax\,\% at random phases; at $B = \WakeBLargeMs$~ms, here and on
Windows, it is checked at one phase only. $\Beff$ exceeds $B$ by \WakeEpollBEffExcessOneMs~$\mu$s
(epoll) and \WakeUringBEffExcessOneMs~$\mu$s (io\_uring) at \WakeBSmallMs~ms, by
\WakeSweepEpollBEffExcessTenMs\ and \WakeSweepUringBEffExcessTenMs~$\mu$s at \WakeBMidMs~ms, and
by \WakeEpollBEffExcessOneHundredMs\ and \WakeUringBEffExcessOneHundredMs~$\mu$s at
\WakeBLargeMs~ms. The default timer slack adds \WakeEpollSlackEffectOneMs~$\mu$s to the epoll
period and leaves io\_uring at \WakeUringBEffOneMsDefaultSlack~$\mu$s; the man page lists
\texttt{epoll\_wait}, not \texttt{io\_\allowbreak{}uring\_\allowbreak{}enter}, among the calls that
slack affects~\cite{manpages2026timerslack}.

\begin{figure}[t]
\centering
% Measured delay against phase, defect arm, B = 1 ms: one panel per backend.
% Data: results/rework/fig_sawtooth.csv, made by
% analysis/fig_sawtooth.py --results results/rework from gap_sweep_posts.csv and
% gap_sweep_summary.csv in the same directory. No cell is bimodal, so there is no kind 2.
% Series 1 epoll (L), 2 io_uring (L), 3 IOCP (W, default timer). Kind 0 is a measured post;
% kind 1 the model d = B_eff - (phase mod B_eff) with the pooled B_eff of the cell.
% One column wide; the caption and label belong in main.tex. Needs pgfplots (compat 1.18).
\pgfplotstableread[col sep=comma]{../results/rework/fig_sawtooth.csv}\WakeSawtoothTable
\begin{tikzpicture}[
    /pgfplots/sawtooth panel/.style={
        scale only axis, width=0.78\columnwidth, height=2.2cm,
        xmin=0, ymin=0, enlarge x limits={upper, value=0.02}, enlarge y limits={upper, value=0.32},
        scaled ticks=false,
        tick label style={font=\scriptsize, /pgf/number format/1000 sep={}},
        label style={font=\footnotesize},
        ylabel={Delay ($\mu$s)},
        every axis y label/.style={at={(axis description cs:-0.15,0.5)}, rotate=90, anchor=south},
        tick align=outside, tick pos=left,
    },
    /pgfplots/sawtooth posts/.style={
        only marks, mark=*, mark size=0.8pt, color=#1,
        mark options={draw opacity=0, fill=#1, fill opacity=0.65},
    },
    /pgfplots/sawtooth model/.style={no marks, thin, black!70, unbounded coords=jump},
    /pgfplots/sawtooth select/.style={restrict expr to domain={10*\thisrow{series}+\thisrow{kind}}{#1}},
    panel label/.style={anchor=north east, font=\scriptsize, inner sep=1.5pt},
]
\definecolor{sawtoothEpoll}{RGB}{0,114,178}
\definecolor{sawtoothUring}{RGB}{213,94,0}
\definecolor{sawtoothIocp}{RGB}{0,158,115}

\begin{axis}[sawtooth panel, name=epoll,
    legend style={font=\scriptsize, at={(0.5,1.04)}, anchor=south, legend columns=2, draw=none,
                  /tikz/every even column/.append style={column sep=8pt}}]
    \addplot[sawtooth model, sawtooth select=11:11] table[x=phase_us, y=delay_us] {\WakeSawtoothTable};
    \addlegendentry{model $B_{\mathrm{eff}} - (\text{phase} \bmod B_{\mathrm{eff}})$}
    \addplot[sawtooth posts=sawtoothEpoll, sawtooth select=10:10]
        table[x=phase_us, y=delay_us] {\WakeSawtoothTable};
    \addlegendentry{measured post}
    \node[panel label] at (rel axis cs:0.99,0.98) {epoll, Linux};
\end{axis}

\begin{axis}[sawtooth panel, name=uring,
    at={(epoll.below south west)}, anchor=north west, yshift=-2pt]
    \addplot[sawtooth model, sawtooth select=21:21] table[x=phase_us, y=delay_us] {\WakeSawtoothTable};
    \addplot[sawtooth posts=sawtoothUring, sawtooth select=20:20]
        table[x=phase_us, y=delay_us] {\WakeSawtoothTable};
    \node[panel label] at (rel axis cs:0.99,0.98) {io\_uring, Linux};
\end{axis}

\begin{axis}[sawtooth panel, name=iocp,
    at={(uring.below south west)}, anchor=north west, yshift=-2pt,
    xlabel={Phase since the worker parked ($\mu$s)}]
    \addplot[sawtooth model, sawtooth select=31:31] table[x=phase_us, y=delay_us] {\WakeSawtoothTable};
    \addplot[sawtooth posts=sawtoothIocp, sawtooth select=30:30]
        table[x=phase_us, y=delay_us] {\WakeSawtoothTable};
    \node[panel label] at (rel axis cs:0.99,0.98) {IOCP, Windows, default timer};
\end{axis}
\end{tikzpicture}

\caption{Delay of posted tasks against their phase in a defective loop with
$B = \WakeBSmallMs$~ms: epoll and io\_uring on Linux, IOCP on Windows with the default timer
resolution. Every second post of each run is shown (\WakeFigPostsShown\ of \WakeSweepCellPosts\
per panel). The line is (\ref{eq:model}) with the measured $\Beff$.}
\label{fig:sawtooth}
\end{figure}

On Windows the wait ends on a timer tick. At the default resolution, bounds of \WakeBSmallMs\
and \WakeBMidMs~ms both give one period of \WakeWinDefaultTickMs~ms, and \WakeBLargeMs~ms gives
\WakeIocpBEffOneHundredMsDefaultTimer~$\mu$s, \WakeIocpBEffOverTicksOneHundredMsDefaultTimer~$\mu$s
more than \WakeIocpTicksOneHundredMsDefaultTimer\ times that measured tick. The largest residual
at this resolution is \WakeIocpAbsResidualPctOneMsDefaultTimer\,\%, so H1 holds with the default
timer, and the largest delay of a
defective loop at bounds of \WakeBSmallMs\ and \WakeBMidMs~ms is
\WakeIocpDefectMaxMsDefaultTimer~ms. With \WakeTimerRequestedMs~ms requested,
bounds of \WakeBMidMs\ and \WakeBLargeMs~ms give periods \WakeSweepIocpBEffExcessTenMsTimerOneMs\
and \WakeIocpBEffExcessOneHundredMsTimerOneMs~$\mu$s longer than $B$, with residuals of
\WakeSweepIocpAbsResidualPctTenMsTimerOneMs\ and \WakeIocpAbsResidualPctOneHundredMsTimerOneMs\,\%.
A \WakeBSmallMs~ms bound with the \WakeTimerRequestedMs~ms timer gives a median period of
\WakeSweepIocpBEffOneMsTimerOneMs~$\mu$s. The sweep then loses control of the phase:
\WakeGridPostsIocpOneMsTimerOneMs\ of \WakeSweepCellPosts\ posts land within \WakeGridTolUs~$\mu$s
of a multiple of \WakeGridUs~$\mu$s (\WakeGridUniformPct\,\% if spread evenly), and
\WakeBandPostsIocpOneMsTimerOneMs\ are delivered after \WakeBandLoUsIocpOneMsTimerOneMs\ to
\WakeBandHiUsIocpOneMsTimerOneMs~$\mu$s whatever their phase. There the model misses by
\WakeSweepIocpAbsResidualPctOneMsTimerOneMs\,\% of $\Beff$ (the matrix cell, at one phase, by
\WakeIocpAbsResidualPctOneMsTimerOneMs\,\%), a failure we do not explain, and H1 fails for the
\WakeTimerRequestedMs~ms timer.

\emph{Detection.} In the matrix the test flagged every defective configuration and passed
every correct one (\WakeDetectConfigsLinux\ per build on Linux, \WakeDetectConfigsWindows\ on
Windows). Defective medians exceed $\Beff/\WakeTestDivisor$ by factors of
\WakeDetectFactorDefectMin\ to \WakeDetectFactorDefectMax, and the largest correct median is
\WakeDetectFactorFixedMax\ of it. $B/\WakeTestDivisor$ gives the same verdicts (factors
\WakeDetectFactorBDefectMin\ to \WakeDetectFactorBDefectMax). But the matrix has one gap per
configuration. The gap sweep runs the same protocol at \WakeSweepSteps\ gaps, and there the test
misses in \WakeSweepMissedTotal\ of \WakeSweepRunsTotal\ runs (\WakeSweepMissedBTotal\ with
$B/\WakeTestDivisor$), \WakeSweepMissedLateTotal\ of them with phases in the last
$\Beff/\WakeTestDivisor$ of the period. For epoll at $B = \WakeBMidMs$~ms and a
\WakeSweepEpollMinMedianGapMsTenMs~ms gap, the defective median is
\WakeSweepEpollMinMedianTenMs~$\mu$s, close to the \WakeBoundEpollFixedTenMs~$\mu$s of the
correct loop at that bound. Random phases shrink this blind spot. Drawing \WakeTestPosts\ posts
at random from each sweep, the largest miss rate is \WakeResampleMissMaxPct\,\% of draws
(\WakeResampleMissBMaxPct\,\% with $B/\WakeTestDivisor$), near the \WakeBinomialMissPct\,\% of
Section~\ref{sec:test}. In the exploratory random-gap run on Linux the test missed
\WakeRandomDefectMissed\ of \WakeRandomDefectTrials\ defective trials
(\WakeRandomDefectMissPct\,\%; \WakeRandomDefectMissedB\ with $B/\WakeTestDivisor$) and flagged
\WakeRandomFixedFlagged\ of \WakeRandomFixedTrials\ correct ones. For Windows, only the
resampling is available.

\emph{Cost of a correct wake.} H2 expected the delay of a correct loop not to depend on $B$. On
Windows this holds, with medians of \WakeBoundIocpFixedBlockingTimerOneMs\ to
\WakeBoundIocpFixedTenMsTimerOneMs~$\mu$s with the \WakeTimerRequestedMs~ms timer and
\WakeBoundIocpFixedBlockingDefaultTimer\ to \WakeBoundIocpFixedOneHundredMsDefaultTimer~$\mu$s
with the default. On Linux it fails: at a \WakeBoundGapMs~ms gap, epoll delivers in
\WakeBoundEpollFixedOneMs~$\mu$s at $B = \WakeBSmallMs$~ms but in \WakeBoundEpollFixedTenMs\ to
\WakeBoundEpollFixedOneHundredMs~$\mu$s with longer bounds or a blocking wait, and io\_uring in
\WakeBoundUringFixedOneMs~$\mu$s against \WakeBoundUringFixedOneHundredMs\ to
\WakeBoundUringFixedTenMs~$\mu$s. Under the \WakeQosUs~$\mu$s latency request, which leaves idle
cores only a polling state~\cite{kernel2026cpuidle}, every median is lower: epoll delivers in
\WakeBoundQosEpollFixedOneMsQosZeroUs~$\mu$s at $B = \WakeBSmallMs$~ms and in
\WakeBoundQosEpollFixedBlockingQosZeroUs\ to \WakeBoundQosEpollFixedTenMsQosZeroUs~$\mu$s
otherwise, and io\_uring in \WakeBoundQosUringFixedOneMsQosZeroUs\ and
\WakeBoundQosUringFixedBlockingQosZeroUs\ to \WakeBoundQosUringFixedTenMsQosZeroUs~$\mu$s. The
dependence on $B$ remains, so the exploratory H5 fails for both backends. With and without the
request, the \WakeBSmallMs~ms bound gives the fastest delivery. We did not record idle-state
residency, and the data do not show why the shortest bound is fastest.

\emph{Threats to validity.} Each system ran on one machine, and the machines differ in
processor as well as operating system, so the contrast between Linux and Windows in H2 may come
from either. All runs use an idle loop without I/O, where the
defect shows. Under load, completions end the wait and hide it, so our delays bound its cost
from above and say nothing about how often it occurs in production. Each cell ran once, so the
variation between runs is unknown. The loop is ours and minimal, and loops with I/O, timers or
several workers may wake differently. The seeded faults are few and hand-written, and kqueue is
not measured. The test assumes that a correct wake costs much less than $B/\WakeTestDivisor$,
which a short bound, a sanitizer build or a deep idle state can break. The random-gap design and
the resampling of the gap sweep were added after a review of results from the earlier
implementation and were not pre-specified.

\section{Conclusion}

A wake that misses a bounded wait does not break functional correctness but costs up to one
wait period in latency, so functional tests pass. Its delay is a sawtooth in the phase of the
post, except for a \WakeBSmallMs~ms bound under a \WakeTimerRequestedMs~ms timer on Windows. A
test that posts a few tasks to an idle worker and compares their median delay with
$\Beff/\WakeTestDivisor$ flagged every seeded defect in our matrix; a fixed gap can hide the
defect, and random gaps shrink that blind spot. On Windows the timer resolution of the process
rounds the period up: at the default resolution a defect delayed a task by up to one period of
about \WakeWinDefaultTickMs~ms at bounds of \WakeBSmallMs\ and \WakeBMidMs~ms; the largest delay
in our runs was \WakeIocpDefectMaxMsDefaultTimer~ms. With a correct wake, the delay on an idle Linux
core still depends on the bound, also under a \WakeQosUs~$\mu$s CPU latency request. We
recommend running the latency test with random gaps alongside functional tests, and controlling
or reporting the idle policy in benchmarks of wake latency.

\bibliographystyle{IEEEtran}
\bibliography{refs}

\end{document}
